\documentclass[11pt, a4paper]{article}

\usepackage[utf8]{inputenc}
\usepackage[T1]{fontenc}
\usepackage{amsmath, amssymb, amsthm, mathtools}
\usepackage{bm}
% indicator function via boldsymbol (no extra package needed)
\usepackage{booktabs}
\usepackage{graphicx}
\usepackage{hyperref}
\usepackage[margin=2.5cm]{geometry}
\usepackage{enumitem}
\usepackage{xcolor}
\usepackage{tcolorbox}
\usepackage{float}

% Theorem environments
\theoremstyle{definition}
\newtheorem{definition}{Definition}[section]
\newtheorem{example}{Example}[section]
\theoremstyle{plain}
\newtheorem{theorem}{Theorem}[section]
\newtheorem{proposition}{Proposition}[section]
\newtheorem{corollary}{Corollary}[section]
\theoremstyle{remark}
\newtheorem{remark}{Remark}[section]

% Colored boxes
\newtcolorbox{keyresult}[1][]{
    colback=blue!5!white,
    colframe=blue!60!black,
    fonttitle=\bfseries,
    title={Key Insight},
    #1
}
\newtcolorbox{warningbox}[1][]{
    colback=red!5!white,
    colframe=red!60!black,
    fonttitle=\bfseries,
    title={Warning},
    #1
}
\newtcolorbox{refbox}[1][]{
    colback=gray!5!white,
    colframe=gray!50!black,
    fonttitle=\bfseries,
    title={Go Deeper},
    #1
}
\newtcolorbox{keyidea}[1][]{
    colback=green!5!white,
    colframe=green!50!black,
    fonttitle=\bfseries,
    title={Core Idea},
    #1
}
\newtcolorbox{prereq}[1][]{
    colback=orange!5!white,
    colframe=orange!60!black,
    fonttitle=\bfseries,
    title={Read Before Coding},
    #1
}
\newtcolorbox{objective}[1][]{
    colback=blue!5!white,
    colframe=blue!60!black,
    fonttitle=\bfseries,
    title={What You Will Build},
    #1
}

% Custom commands
\newcommand{\VaR}{\operatorname{VaR}}
\newcommand{\ES}{\operatorname{ES}}
\newcommand{\CVaR}{\operatorname{CVaR}}
\newcommand{\EV}{\mathbb{E}}
\newcommand{\PV}{\mathbb{P}}
\newcommand{\RV}{\mathbb{R}}
\newcommand{\ind}{\boldsymbol{1}}

\title{%
    \textbf{Copula Dependency Modeling} \\[0.5em]
    \large Beyond Correlation: Modeling How Risks Move Together
}
\author{Andr\'es Gonz\'alez Ortega}
\date{March 2026}

\begin{document}
\maketitle

\begin{abstract}
An applied treatment of copula theory for portfolio risk modeling.
We begin with the concrete failure of Pearson correlation during the 2008 crisis,
then develop the mathematical framework: Sklar's theorem, elliptical copulas, Archimedean copulas,
and vine constructions for high dimensions.
A complete workflow --- GARCH filtering, probability integral transform, copula fitting, and joint simulation ---
is described and applied to a three-asset portfolio.
The main result: choosing the wrong copula (Gaussian) underestimates portfolio tail risk
by 20--40\% compared to copulas that capture joint crashes.
We assume strong probability/statistics background and focus on the financial risk application.
\end{abstract}

\tableofcontents
\newpage

%% ============================================================================
\section{Why Correlation Fails}
\label{sec:why-copulas}
%% ============================================================================

\begin{example}[Correlation breakdown in 2008]\label{ex:corr-breakdown}
From 2005 to mid-2007, the rolling 252-day correlation between US equity indices
and mortgage-backed securities (MBS) was approximately $\rho \approx 0.3$.
Portfolio managers used this number to ``diversify'':
if equities fall, MBS won't fall as much, and vice versa.

During the 2008 financial crisis, the correlation spiked above $\rho \approx 0.8$.
The ``diversification'' vanished precisely when it was needed most.
A portfolio optimized under $\rho = 0.3$ suffered losses far larger than any
model based on constant correlation had predicted.
\end{example}

This is not an isolated anecdote.
Correlation fails as a dependence measure for at least three fundamental reasons:

\begin{enumerate}
    \item \textbf{Linearity.}
    Pearson correlation measures only \emph{linear} association.
    If $X \sim N(0,1)$ and $Y = X^2$, then $X$ and $Y$ are perfectly dependent ($Y$ is a deterministic function of $X$),
    yet $\rho(X, Y) = 0$.
    Correlation misses all non-linear dependence.

    \item \textbf{Marginal dependence.}
    Correlation depends on the marginal distributions, not just the dependence structure.
    The same copula (same dependence) paired with different marginals can produce different correlations.
    This means correlation is not an intrinsic property of dependence.

    \item \textbf{Tail blindness.}
    Two joint distributions can have identical correlation $\rho$ but radically different
    probabilities of joint extreme events.
    A Gaussian copula and a Student-$t$ copula with the same $\rho$ agree in the center
    but diverge dramatically in the tails.
\end{enumerate}

\begin{example}[The Gaussian copula catastrophe]\label{ex:gaussian-copula}
In the early 2000s, the Gaussian copula (popularized by David Li's 2000 paper)
became the market standard for pricing Collateralized Debt Obligations (CDOs).
The model assumed that joint defaults were governed by a multivariate normal dependence structure.

The fatal flaw: the Gaussian copula has \emph{zero tail dependence}.
It assigns negligible probability to many assets defaulting simultaneously.
When the housing market collapsed in 2007--2008, widespread simultaneous defaults occurred ---
events the Gaussian copula model deemed essentially impossible.
Senior CDO tranches rated AAA suffered catastrophic losses.

This was not a failure of copula theory.
It was a failure of choosing the \emph{wrong} copula: one that structurally cannot capture joint extremes.
\end{example}

\begin{warningbox}
Correlation is a single number.
Dependence is a \emph{function} --- a full description of how the joint distribution relates to the marginals.
Compressing it into one number discards critical information,
especially about the tails where risk management matters most.
\end{warningbox}

%% ============================================================================
\section{Sklar's Theorem: The Fundamental Decomposition}
\label{sec:sklar}
%% ============================================================================

\begin{definition}[Copula]
A \textbf{$d$-dimensional copula} is a function $C : [0,1]^d \to [0,1]$ such that:
\begin{enumerate}
    \item $C$ is grounded: $C(u_1, \ldots, u_d) = 0$ if any $u_i = 0$.
    \item $C$ has uniform marginals: $C(1, \ldots, 1, u_i, 1, \ldots, 1) = u_i$ for all $i$.
    \item $C$ is $d$-increasing (the copula analogue of a CDF being non-decreasing).
\end{enumerate}
\end{definition}

\begin{theorem}[Sklar, 1959]\label{thm:sklar}
Let $F$ be a $d$-dimensional joint distribution function with marginals $F_1, \ldots, F_d$.
Then there exists a copula $C$ such that for all $(x_1, \ldots, x_d) \in \overline{\RV}^d$:
\[
    F(x_1, \ldots, x_d) = C\bigl(F_1(x_1), \ldots, F_d(x_d)\bigr)
\]
If the marginals $F_1, \ldots, F_d$ are continuous, then $C$ is \textbf{unique}.

Conversely, if $C$ is a copula and $F_1, \ldots, F_d$ are univariate distribution functions,
then the function $F$ defined above is a valid joint distribution with marginals $F_1, \ldots, F_d$.
\end{theorem}

\begin{keyresult}
Sklar's theorem provides a \textbf{complete separation} of two modeling tasks:
\begin{enumerate}
    \item \textbf{Marginals}: model each risk factor's distribution individually
          (e.g., GARCH with skewed-$t$ innovations, or GARCH + EVT from Project~05).
    \item \textbf{Dependence}: model how the risk factors move together by choosing a copula $C$.
\end{enumerate}
The two choices are \textbf{independent}.
You can pair any set of marginals with any copula.
This modularity is the main practical advantage of the copula approach.
\end{keyresult}

The copula itself can be extracted from a known joint distribution.
If $F$ has continuous marginals, then:
\[
    C(u_1, \ldots, u_d) = F\bigl(F_1^{-1}(u_1), \ldots, F_d^{-1}(u_d)\bigr)
\]
This is how the Gaussian and $t$-copulas are derived: extract the copula from the multivariate
normal or multivariate $t$ distribution.

%% ============================================================================
\section{Elliptical Copulas}
\label{sec:elliptical}
%% ============================================================================

Elliptical copulas are derived from elliptical distributions (multivariate normal, multivariate $t$, etc.).
They inherit the symmetry of the parent distribution: the dependence structure
is the same in all directions.

\subsection{Gaussian Copula}

\begin{definition}[Gaussian copula]
The $d$-dimensional Gaussian copula with correlation matrix $\bm{R}$ is:
\[
    C_{\bm{R}}^{\text{Ga}}(u_1, \ldots, u_d)
    = \Phi_{\bm{R}}\bigl(\Phi^{-1}(u_1), \ldots, \Phi^{-1}(u_d)\bigr)
\]
where $\Phi$ is the standard normal CDF and $\Phi_{\bm{R}}$ is the CDF of the
$d$-dimensional normal with mean $\bm{0}$ and correlation matrix $\bm{R}$.
\end{definition}

The Gaussian copula is parameterized by the $d(d-1)/2$ off-diagonal entries of $\bm{R}$.
It captures the \emph{rank} dependence structure of the multivariate normal,
without imposing normal marginals.

\begin{proposition}[Tail dependence of Gaussian copula]
For any $\rho \in (-1, 1)$, the Gaussian copula has
\[
    \lambda_L = \lambda_U = 0
\]
The Gaussian copula is \textbf{asymptotically independent} in the tails:
the probability of joint extreme events vanishes faster than the probability of individual extremes.
\end{proposition}

This is the mathematical reason the Gaussian copula failed for CDO pricing:
it structurally cannot generate the joint extreme events that matter for risk.

\subsection{Student-$t$ Copula}

\begin{definition}[Student-$t$ copula]
The $d$-dimensional $t$-copula with correlation matrix $\bm{R}$ and $\nu$ degrees of freedom is:
\[
    C_{\bm{R}, \nu}^{t}(u_1, \ldots, u_d)
    = t_{\bm{R}, \nu}\bigl(t_\nu^{-1}(u_1), \ldots, t_\nu^{-1}(u_d)\bigr)
\]
where $t_\nu$ is the univariate $t$-CDF and $t_{\bm{R}, \nu}$ is the CDF of the
$d$-dimensional $t$-distribution with correlation $\bm{R}$ and $\nu$ degrees of freedom.
\end{definition}

The $t$-copula has one extra parameter compared to the Gaussian: the degrees of freedom $\nu$.
As $\nu \to \infty$, it converges to the Gaussian copula.

\begin{proposition}[Tail dependence of $t$-copula]
For the bivariate $t$-copula with correlation $\rho$ and $\nu$ degrees of freedom:
\[
    \lambda_L = \lambda_U = 2 \, t_{\nu+1}\!\left(-\sqrt{\nu + 1}\sqrt{\frac{1-\rho}{1+\rho}}\right)
\]
where $t_{\nu+1}$ is the CDF of the univariate $t$-distribution with $\nu + 1$ degrees of freedom.
\end{proposition}

\begin{example}[Tail dependence values]
For $\rho = 0.5$:
\begin{itemize}
    \item $\nu = 4$: $\lambda_L = \lambda_U \approx 0.18$ (substantial tail dependence).
    \item $\nu = 10$: $\lambda_L = \lambda_U \approx 0.07$ (moderate).
    \item $\nu = 30$: $\lambda_L = \lambda_U \approx 0.01$ (weak, approaching Gaussian).
    \item $\nu \to \infty$: $\lambda_L = \lambda_U = 0$ (Gaussian limit).
\end{itemize}
Lower $\nu$ means heavier joint tails.
For equity data, typical estimates are $\nu \approx 5$--$10$.
\end{example}

\subsection{Comparison}

\begin{table}[H]
\centering
\begin{tabular}{@{}lcccc@{}}
\toprule
\textbf{Property} & \textbf{Gaussian} & \textbf{Student-$t$} \\
\midrule
Parameters & $\bm{R}$ ($d(d-1)/2$ entries) & $\bm{R}$ + $\nu$ \\
Tail dependence & $\lambda_L = \lambda_U = 0$ & $\lambda_L = \lambda_U > 0$ \\
Symmetry & Symmetric & Symmetric \\
Joint extremes & Cannot capture & Captures \\
Nests Gaussian? & --- & Yes ($\nu \to \infty$) \\
Typical use & Benchmark / comparison & Default choice for multi-asset risk \\
\bottomrule
\end{tabular}
\caption{Gaussian vs.\ Student-$t$ copula.
The $t$-copula is strictly more flexible and captures the joint tail behavior
that matters for risk.}
\end{table}

\begin{keyresult}
The tail dependence coefficient $\lambda$ has a direct risk interpretation:
$\lambda_L = \PV(U_2 \leq \epsilon \mid U_1 \leq \epsilon)$ as $\epsilon \to 0$,
where $U_1, U_2$ are the copula-uniform variables.
In portfolio terms: \emph{given that asset 1 has an extreme loss,
what is the probability that asset 2 simultaneously has an extreme loss?}
The Gaussian copula answers: zero.
The $t$-copula gives a realistic, positive number.
\end{keyresult}

%% ============================================================================
\section{Archimedean Copulas}
\label{sec:archimedean}
%% ============================================================================

Archimedean copulas offer a different construction from elliptical copulas.
They are built from a \emph{generator function} and can capture asymmetric tail dependence
(e.g., lower tail dependence without upper, or vice versa).

\subsection{The Generator Construction}

\begin{definition}[Archimedean copula]
Let $\varphi : [0,1] \to [0, \infty]$ be a continuous, strictly decreasing, convex function
with $\varphi(1) = 0$ (the \textbf{generator}).
Let $\varphi^{-1}$ be its inverse (the \textbf{pseudo-inverse} if $\varphi(0) < \infty$).
The bivariate Archimedean copula is:
\[
    C(u_1, u_2) = \varphi^{-1}\bigl(\varphi(u_1) + \varphi(u_2)\bigr)
\]
\end{definition}

The generator function encodes all the dependence structure.
Different generators give different copula families with distinct tail properties.

\subsection{The Three Main Families}

\begin{table}[H]
\centering
\begin{tabular}{@{}lllcc@{}}
\toprule
\textbf{Family} & \textbf{Generator} $\varphi(t)$ & \textbf{Parameter} & $\lambda_L$ & $\lambda_U$ \\
\midrule
Clayton & $\frac{1}{\theta}(t^{-\theta} - 1)$ & $\theta > 0$ & $2^{-1/\theta}$ & 0 \\[6pt]
Gumbel & $(-\ln t)^\theta$ & $\theta \geq 1$ & 0 & $2 - 2^{1/\theta}$ \\[6pt]
Frank & $-\ln\!\left(\frac{e^{-\theta t} - 1}{e^{-\theta} - 1}\right)$ & $\theta \neq 0$ & 0 & 0 \\
\bottomrule
\end{tabular}
\caption{The three main Archimedean copula families and their tail dependence properties.}
\end{table}

\textbf{Clayton} ($\theta > 0$):
\[
    C_\theta^{\text{Cl}}(u_1, u_2) = \left(u_1^{-\theta} + u_2^{-\theta} - 1\right)^{-1/\theta}
\]
\begin{itemize}
    \item Strong \emph{lower} tail dependence, zero upper tail dependence.
    \item As $\theta \to 0^+$, approaches independence; as $\theta \to \infty$, approaches comonotonicity.
    \item \textbf{Use case}: modeling joint crashes. When both assets are in their lower tails simultaneously, Clayton assigns higher probability than the Gaussian or Frank copulas.
\end{itemize}

\textbf{Gumbel} ($\theta \geq 1$):
\[
    C_\theta^{\text{Gu}}(u_1, u_2) = \exp\!\left\{-\left[(-\ln u_1)^\theta + (-\ln u_2)^\theta\right]^{1/\theta}\right\}
\]
\begin{itemize}
    \item Strong \emph{upper} tail dependence, zero lower tail dependence.
    \item At $\theta = 1$, reduces to the independence copula $\Pi(u_1, u_2) = u_1 u_2$.
    \item \textbf{Use case}: modeling joint booms or simultaneous extreme gains (less common in risk, but relevant for insurance).
\end{itemize}

\textbf{Frank} ($\theta \neq 0$):
\[
    C_\theta^{\text{Fr}}(u_1, u_2) = -\frac{1}{\theta}\ln\!\left(1 + \frac{(e^{-\theta u_1} - 1)(e^{-\theta u_2} - 1)}{e^{-\theta} - 1}\right)
\]
\begin{itemize}
    \item \emph{No} tail dependence in either direction.
    \item Symmetric dependence.
    \item \textbf{Use case}: when you need a copula with no tail dependence but want flexibility
          beyond the Gaussian (Frank allows a wider range of concordance measures).
\end{itemize}

\subsection{Worked Example: Clayton Copula for Joint Crashes}

\begin{example}[Fitting Clayton to an equity pair]\label{ex:clayton}
Consider daily losses of two assets: a US equity ETF (SPY) and a European equity ETF (EZU), 2000--2024.
After GARCH filtering and PIT (described in Section~\ref{sec:workflow}),
we have approximately 6{,}000 pairs $(u_1, u_2) \in [0,1]^2$.

Fitting the Clayton copula by maximum likelihood:
$\hat\theta = 0.85$.

\textbf{Implied lower tail dependence}:
\[
    \hat\lambda_L = 2^{-1/\hat\theta} = 2^{-1/0.85} \approx 0.44
\]

Interpretation: given that SPY is experiencing an extreme loss
(say, in its worst 1\% of days), there is a 44\% probability that EZU
is \emph{simultaneously} in its worst 1\%.
Under the Gaussian copula (regardless of correlation), this probability would approach 0\%.

The Clayton copula captures what we see in crises:
when one equity market crashes, others tend to crash simultaneously.
The ``diversification benefit'' of holding both US and European equities is much smaller
in the tail than correlation alone suggests.
\end{example}

%% ============================================================================
\section{Vine Copulas}
\label{sec:vine}
%% ============================================================================

For $d = 2$, any of the copulas above work directly.
For $d > 2$, the situation is more complex:
Archimedean copulas impose the same dependence between every pair (exchangeability),
which is unrealistic.
Elliptical copulas handle multiple dimensions but cannot capture asymmetric tail dependence.

\textbf{Vine copulas} solve this by decomposing a $d$-dimensional copula
into $d(d-1)/2$ \emph{bivariate} copulas, organized in a tree structure.
Each bivariate building block (``pair copula'') can be from a different family.

\subsection{Pair-Copula Construction}

The idea relies on the decomposition of a $d$-dimensional density into conditional bivariate densities.
For $d = 3$ with variables $(X_1, X_2, X_3)$:
\begin{align*}
    f(x_1, x_2, x_3)
    &= f_1(x_1) \cdot f_{2|1}(x_2 \mid x_1) \cdot f_{3|12}(x_3 \mid x_1, x_2) \\
    &= f_1(x_1) \cdot f_2(x_2) \cdot c_{12}(F_1(x_1), F_2(x_2)) \\
    &\quad \times f_3(x_3) \cdot c_{13|2}(F_{1|2}(x_1 \mid x_2), F_{3|2}(x_3 \mid x_2))
\end{align*}
where $c_{ij}$ and $c_{ij|k}$ are bivariate copula densities (conditional and unconditional).

The $\binom{d}{2} = d(d-1)/2$ bivariate copulas are organized in a sequence of trees.
For $d = 5$, we need $5 \times 4 / 2 = 10$ pair copulas.

\subsection{Three Vine Structures}

\begin{table}[H]
\centering
\begin{tabular}{@{}lp{6.5cm}l@{}}
\toprule
\textbf{Type} & \textbf{Structure} & \textbf{Best for} \\
\midrule
C-vine & Star: one ``central'' variable connected to all others in tree~1 & One dominant risk factor \\
D-vine & Path: variables in a sequence & Natural ordering (e.g., time or geography) \\
R-vine & General: any valid tree sequence & Most flexible; selected by AIC/BIC \\
\bottomrule
\end{tabular}
\caption{Vine copula structures. The R-vine nests both C-vine and D-vine as special cases.}
\end{table}

\begin{keyidea}
\textbf{Why vine copulas matter for risk.}
A 10-asset portfolio has 45 pairs.
With a vine copula, each pair can have its own copula family:
\begin{itemize}
    \item SPY--EZU: Clayton (strong lower tail dependence, joint equity crashes).
    \item SPY--TLT: Frank (little tail dependence between equities and Treasuries).
    \item GLD--SPY: Gumbel survival copula (dependence mainly during extreme sell-offs).
\end{itemize}
This level of granularity is impossible with a single multivariate copula.
\end{keyidea}

\begin{refbox}
\begin{itemize}
    \item Aas, K., Czado, C., Frigessi, A.\ \& Bakken, H.\ (2009).
          ``Pair-copula constructions of multiple dependence.''
          \emph{Insurance: Mathematics and Economics}, 44(2), 182--198.
          The key applied paper for vine copulas.
    \item Joe, H.\ (2014). \emph{Dependence Modeling with Copulas}.
          Chapman \& Hall/CRC.
          Comprehensive treatment of vine copulas and pair-copula constructions.
    \item Czado, C.\ (2019). \emph{Analyzing Dependent Data with Vine Copulas}.
          Springer. Modern monograph with R examples.
\end{itemize}
\end{refbox}

%% ============================================================================
\section{The Workflow: GARCH Filtering + PIT + Copula}
\label{sec:workflow}
%% ============================================================================

Copulas operate on \emph{uniform} random variables.
Raw financial returns are neither uniform nor independent.
The standard workflow transforms returns into a form suitable for copula fitting.

\subsection{Step 1: GARCH Filtering}

For each asset $i = 1, \ldots, d$, fit a GARCH-type model (e.g., GJR-GARCH with skewed-$t$ innovations):
\[
    r_{i,t} = \hat\mu_{i,t} + \hat\sigma_{i,t} \, z_{i,t}
\]
Extract the standardized residuals:
\[
    \hat{z}_{i,t} = \frac{r_{i,t} - \hat\mu_{i,t}}{\hat\sigma_{i,t}}
\]

\textbf{Why}: this removes the serial dependence (volatility clustering) from each return series.
The residuals $\hat{z}_{i,t}$ should be approximately i.i.d.
(Check with Ljung-Box tests on $\hat{z}_{i,t}$ and $\hat{z}_{i,t}^2$.)

\subsection{Step 2: Probability Integral Transform (PIT)}

Map each standardized residual to the unit interval using the (estimated) CDF of the innovation distribution:
\[
    \hat{u}_{i,t} = \hat{F}_{z_i}(\hat{z}_{i,t})
\]

Two approaches:
\begin{itemize}
    \item \textbf{Parametric PIT}: use the fitted innovation CDF (e.g., skewed-$t$ with estimated $\nu, \lambda$).
    \item \textbf{Semi-parametric PIT}: use the empirical CDF of the residuals (the ``rank transform''):
          $\hat{u}_{i,t} = \text{rank}(\hat{z}_{i,t}) / (T+1)$.
          This avoids misspecification of the marginal distribution.
\end{itemize}

\textbf{Result}: if the marginal model is correct, $\hat{u}_{i,t} \sim \text{Uniform}(0,1)$ independently over time.
The $d$-dimensional vectors $(\hat{u}_{1,t}, \ldots, \hat{u}_{d,t})$ are the \textbf{pseudo-observations}
on which we fit the copula.

\begin{warningbox}
The PIT is the critical link between marginal and copula modeling.
If the marginal model is wrong (e.g., assuming normal innovations when the data is heavy-tailed),
the transformed values $\hat{u}_{i,t}$ will \emph{not} be uniform,
and the copula fit will be contaminated.
Always check uniformity of the PIT values (histogram, Kolmogorov-Smirnov test).
\end{warningbox}

\subsection{Step 3: Copula Fitting}

Fit a copula to the pseudo-observations $\{(\hat{u}_{1,t}, \ldots, \hat{u}_{d,t})\}_{t=1}^T$.

Methods:
\begin{enumerate}
    \item \textbf{Maximum likelihood}: maximize the copula log-likelihood
    $\sum_{t=1}^T \ln c(\hat{u}_{1,t}, \ldots, \hat{u}_{d,t}; \bm{\theta})$
    where $c$ is the copula density.
    \item \textbf{Inference functions for margins (IFM)}: a two-step approach ---
    first estimate marginal parameters, then copula parameters.
    Computationally simpler; asymptotically efficient under correct specification.
    \item \textbf{Model selection}: compare candidate copulas (Gaussian, $t$, Clayton, Gumbel, vine)
    using AIC, BIC, or Vuong's test.
\end{enumerate}

\subsection{Step 4: Joint Simulation}

To generate $N$ joint return scenarios:
\begin{enumerate}
    \item Sample $N$ vectors $(\tilde{u}_1, \ldots, \tilde{u}_d)$ from the fitted copula.
    \item Invert the PIT: $\tilde{z}_i = \hat{F}_{z_i}^{-1}(\tilde{u}_i)$ to get standardized residuals.
    \item Invert the GARCH: $\tilde{r}_{i} = \hat\mu_{i, T+1} + \hat\sigma_{i, T+1} \cdot \tilde{z}_i$ to get return scenarios.
    \item Compute portfolio loss: $\tilde{L} = -\bm{w}^\top \tilde{\bm{r}}$.
\end{enumerate}

The result is a sample of $N$ portfolio losses that respects:
\begin{itemize}
    \item Each asset's conditional volatility (from GARCH).
    \item Each asset's marginal tail behavior (from the innovation distribution or EVT).
    \item The joint dependence structure (from the copula), including tail dependence.
\end{itemize}

\begin{keyidea}
The four-step pipeline (GARCH $\to$ PIT $\to$ Copula $\to$ Simulation) is the \textbf{standard approach}
in modern multivariate risk modeling.
It cleanly separates three concerns:
volatility dynamics (Step~1), marginal tail shape (Step~2), and dependence structure (Step~3).
\end{keyidea}

%% ============================================================================
\section{Portfolio Risk Under Different Copulas}
\label{sec:portfolio}
%% ============================================================================

\begin{example}[Three-asset portfolio simulation]\label{ex:portfolio}
Consider a portfolio with equal weights ($w_i = 1/3$) in three assets:
\begin{itemize}
    \item SPY (US equities): annualized vol $\approx 18\%$, $\hat\xi_{\text{GPD}} \approx 0.25$.
    \item EZU (European equities): annualized vol $\approx 22\%$, $\hat\xi_{\text{GPD}} \approx 0.28$.
    \item HYG (US high-yield bonds): annualized vol $\approx 10\%$, $\hat\xi_{\text{GPD}} \approx 0.22$.
\end{itemize}

After GARCH filtering and PIT, we fit three copula models to the pseudo-observations
and simulate $N = 100{,}000$ joint scenarios from each.

\medskip
\textbf{Copula specifications}:
\begin{enumerate}
    \item \textbf{Gaussian copula}: correlation matrix estimated from the pseudo-observations.
    Pairwise correlations: $\hat\rho_{12} = 0.62$, $\hat\rho_{13} = 0.55$, $\hat\rho_{23} = 0.48$.

    \item \textbf{Student-$t$ copula}: same correlation matrix, $\hat\nu = 6$ (estimated by MLE).

    \item \textbf{Clayton copula}: $\hat\theta = 1.2$ (estimated from the lower tail concordance).
\end{enumerate}

\medskip
\textbf{Results} (daily portfolio loss, as \% of portfolio value):

\begin{center}
\begin{tabular}{@{}lccc@{}}
\toprule
\textbf{Risk measure} & \textbf{Gaussian} & \textbf{Student-$t$} & \textbf{Clayton} \\
\midrule
$\VaR_{0.95}$ & 1.82\% & 1.89\% & 1.91\% \\
$\VaR_{0.99}$ & 2.65\% & 3.01\% & 3.15\% \\
$\VaR_{0.999}$ & 3.52\% & 4.38\% & 4.71\% \\
$\ES_{0.99}$ & 3.18\% & 3.92\% & 4.24\% \\
\bottomrule
\end{tabular}
\end{center}

\medskip
\textbf{Interpretation}:
\begin{itemize}
    \item At 95\%, the three copulas give similar results.
          In the center of the distribution, the copula choice barely matters.
    \item At 99\%, the $t$-copula gives 14\% higher VaR than Gaussian; Clayton gives 19\% higher.
    \item At 99.9\%, the gap widens: the $t$-copula exceeds Gaussian by 24\%; Clayton by 34\%.
    \item For $\ES_{0.99}$, Clayton exceeds Gaussian by 33\%.
          This is the ``missing'' tail risk that the Gaussian copula systematically ignores.
\end{itemize}

\textbf{Why}: during joint crashes, all three assets fall together.
The Gaussian copula treats extreme co-movements as negligibly rare.
The $t$ and Clayton copulas assign realistic probability to these joint events.
The portfolio loss in a joint crash is much larger than the sum of individual VaRs
would suggest under Gaussian dependence.
\end{example}

\begin{keyresult}
For a \$100M portfolio, the difference between Gaussian and Clayton copulas at 99.9\%
is approximately \$1.2M in VaR and \$1.1M in ES.
This is not a theoretical curiosity --- it directly affects capital reserves.
Using the wrong copula means holding insufficient capital for joint tail events.
\end{keyresult}

\begin{warningbox}
The Clayton copula used here has \emph{exchangeable} dependence:
it assumes the same lower tail dependence between every pair.
In practice, the SPY--EZU pair likely has stronger lower tail dependence than SPY--HYG.
For a more realistic model, use a \textbf{vine copula} (Section~\ref{sec:vine})
where each pair has its own copula family and parameters.
\end{warningbox}

%% ============================================================================
\section{References}
\label{sec:references}
%% ============================================================================

\begin{enumerate}[label={[\arabic*]}, leftmargin=2cm]
    \item \textbf{Nelsen, R.B.}\ (2006).
          \emph{An Introduction to Copulas}. 2nd ed.
          Springer.
          The standard textbook.
          Chapters 2--4 cover Sklar's theorem, Archimedean copulas, and measures of dependence.

    \item \textbf{Joe, H.}\ (2014).
          \emph{Dependence Modeling with Copulas}.
          Chapman \& Hall/CRC.
          Advanced treatment; the definitive reference on vine copulas and pair-copula constructions.

    \item \textbf{Embrechts, P., McNeil, A.J.\ \& Straumann, D.}\ (2002).
          ``Correlation and dependence in risk management: properties and pitfalls.''
          In \emph{Risk Management: Value at Risk and Beyond}, ed.\ M.\ Dempster, pp.\ 176--223.
          Cambridge University Press.
          The paper that made the case for copulas in risk management by demonstrating
          the failures of correlation.

    \item \textbf{Aas, K., Czado, C., Frigessi, A.\ \& Bakken, H.}\ (2009).
          ``Pair-copula constructions of multiple dependence.''
          \emph{Insurance: Mathematics and Economics}, 44(2), 182--198.
          The key applied paper for vine copulas; includes algorithms and financial examples.

    \item \textbf{McNeil, A.J., Frey, R.\ \& Embrechts, P.}\ (2015).
          \emph{Quantitative Risk Management: Concepts, Techniques and Tools}. 2nd ed.
          Princeton University Press.
          Chapter 7 covers copulas in the context of quantitative risk management.

    \item \textbf{Czado, C.}\ (2019).
          \emph{Analyzing Dependent Data with Vine Copulas}.
          Springer.
          Modern monograph with detailed algorithms and R code examples.

    \item \textbf{Li, D.X.}\ (2000).
          ``On default correlation: a copula function approach.''
          \emph{Journal of Fixed Income}, 9(4), 43--54.
          The paper that introduced the Gaussian copula to credit derivatives.
          A cautionary tale.

    \item \textbf{Salmon, F.}\ (2009).
          ``Recipe for disaster: the formula that killed Wall Street.''
          \emph{Wired Magazine}, February 2009.
          Accessible account of the Gaussian copula's role in the financial crisis.
\end{enumerate}

\vfill
\hrule
\smallskip
\noindent\textit{Andr\'es Gonz\'alez Ortega --- Risk Analyst Project} \hfill \textit{March 2026}

\end{document}
